% \section{Applications}

% \subsection{Anomaly Detection}

% We further evaluate SHIFT's robustness and dense localization capabilities by performing zero-shot anomaly segmentation (AS) on the MVTech AD \cite{bergmann_mvtech} and VisA \cite{zou2022_visa_dataset} datasets for industrial anomaly detection (AD).

% For SHIFT, we derive the anomaly map by upsampling the continuous heatmap produced by the learned linear segmentation head. As for all text-steerable methods, we use an ensemble of ten anomaly prompts (e.g., \texttt{``the anomaly in the <object>.''}) and average the resulting heatmaps across prompts.
% % For text-steerable segmentation method such as SHIFT, we apply an ensemble of ten anomaly-related prompts (e.g., \texttt{``the anomaly in the <object>.''}) and average the resulting continuous heatmaps across prompts. 
% Following standard practice in the AD literature \cite{WinCLIP_Jeong23CVPR}, we measure the pixel-wise \textit{Area under the Receiver Operating Characteristics} ($\text{ROC}_P$), the \textit{Per-Region-Overlap} (PRO), and threshold-optimal $F_1$ score ($F_1^{\max}$). Performance of existing training- and adaptation-free AS baselines is stated as reported in the original publication. 

% Despite its general-purpose nature, SHIFT is competitive with dedicated AS approaches. As shown in Table~\ref{tab:mvtec_visa}, it reaches 79.4 PRO on MVTech AD, substantially improving over prompt-based non-AD baselines (e.g., SAM3 at 54.5) and closing much of the gap to specialist methods (e.g., FADE at 84.5). Similar trends appear on VisA, indicating robust transfer to a harder, more diverse inspection setting. Overall, these results suggest that SHIFT’s text-driven post-training unlocks dense visual capabilities already present in the frozen ViT, yields strong out-of-domain generalization to industrial AD tasks without any task-specific training.\todo{use final numbers}

% \begin{table*}[t]
\centering
\footnotesize
\setlength{\tabcolsep}{4pt}
% \renewcommand{\arraystretch}{1.15}
\begin{threeparttable}
\caption{Zero-shot segmentation performance on industrial AD data.}
\label{tab:mvtec_visa}
\begin{tabular}{l l c
                S[table-format=2.1] S[table-format=2.1] S[table-format=2.1]
                S[table-format=2.1] S[table-format=2.1] S[table-format=2.1]}
\toprule
\multicolumn{3}{c}{\textbf{Method}} &
\multicolumn{3}{c}{\textbf{MVTec AD}} &
\multicolumn{3}{c}{\textbf{VisA}} \\
\cmidrule(lr){1-3}\cmidrule(lr){4-6}\cmidrule(lr){7-9}
Name & Ref. & AD &
{$\text{ROC}_P$} & {PRO} & {$F_1^{\max}$} &
{$\text{ROC}_P$} & {PRO} & {$F_1^{\max}$} \\
\midrule
% TransMM  & \cite{Chefer_2021_transmm} & \xmark & 57.5    & 21.9    & 12.1   & 49.4   & 10.2   & 3.1    \\
MaskCLIP & \cite{zhou2022_maskclip} & \xmark & 63.7    & 40.5    & 18.5   & 60.9   & 27.3   & 7.3    \\
CLIPseg  & \cite{lueddecke22_clipseg} & \xmark & 69.0    & 34.6    & 12.5   & 89.5   & 62.4   & 13.9   \\
SAM3     & \cite{carion2025_sam3} & \xmark & 79.9    & 54.5    & 24.1   & 89.8   & 65.9   & 15.5   \\
\addlinespace
\color{gray}WinCLIP  & \color{gray}\cite{WinCLIP_Jeong23CVPR} & \color{gray}\cmark & \color{gray}85.1    & \color{gray}64.6    & \color{gray}31.7   & \color{gray}79.6   & \color{gray}59.8   & \color{gray}14.8   \\
% AnoCLIP  & \cite{deng2024_anoclip} & \cmark & 86.6    & 70.4    & 30.1   & 85.9   & 65.1   & 14.7   \\
% SDP      & \cite{chen2023_clipad} & \cmark & 87.5    & 75.3    & 35.9   & 88.1   & 68.5   & 17.0   \\
% SAA+     & \cite{cao_2023_saa} & \cmark & 73.2    & 42.8    & \underline{37.8}   & 74.0   & 36.8   & \textbf{27.1}   \\
\color{gray} DIVAD    & \color{gray}\cite{hicsonmez2026training}& \color{gray} \cmark & \color{gray}88.0    & \color{gray}72.5    & \color{gray}\underline{35.5}   & \color{gray}\textbf{93.4}   & \color{gray}\underline{78.2}   & \color{gray}\textbf{24.0}   \\
\color{gray}FADE     & \color{gray}\cite{li2024_fade} & \color{gray}\cmark & \color{gray}\textbf{89.6}    & \color{gray}\textbf{84.5}    & \color{gray}\textbf{39.8}   & \color{gray}\underline{91.5}   & \color{gray}\textbf{79.3}   & \color{gray}16.7   \\
\addlinespace
\rowcolor{\methodcolor}\textbf{SHIFT} & \textbf{Ours} & \xmark & \underline{88.1} & \underline{79.4} & 32.6 & \underline{91.5} & 77.5 & \underline{16.9} \\
\bottomrule
\end{tabular}

\begin{tablenotes}[flushleft]\footnotesize
\item AD indicates dedicated anomaly detection methods.
% \item \textdagger{} indicates results reported in one of \cite{WinCLIP_Jeong23CVPR, hicsonmez2026training,chen2023_clipad, deng2024_anoclip}. %mark with \textsuperscript{\textdagger}
\end{tablenotes}
\end{threeparttable}
\end{table*}
